% ABSTRACT - Auto-updated section
% Last generated: 2026-01-15T05:34:34.872Z
% Source: Ollama deepseek-r1:32b

Decentralized Autonomous Organizations (DAOs) represent a paradigm shift in organizational governance, yet their design space remains poorly understood. We present a comprehensive multi-agent simulation framework for analyzing governance mechanisms in DAOs, enabling systematic study of how design parameters affect collective decision-making outcomes. Our framework models heterogeneous agent populations with varying participation strategies, token distributions, and delegation behaviors, supporting multiple voting mechanisms including token-weighted, quadratic, and conviction voting.

Through Monte Carlo simulations (21402 runs across 10 experiment sets), we investigate downstream governance dynamics in DAOs: how proposal pipelines affect decision latency and abandonment, how treasury policy shapes resilience, and how inter-DAO cooperation alters success and resource flows.

We test hypotheses about pipeline design, treasury resilience, and cross-DAO coordination, and report results where experiments have completed. The paper is a living document; claims are updated as additional runs are incorporated. We release our simulation framework as open-source software to enable reproducible governance research and evidence-based mechanism design.

% === AUTO-GENERATION METADATA ===
% This section can be regenerated by running:
%   npm run paper:update -- --section abstract
% Using results from: {{RESULTS_DIR}}
